% GENERATED FILE. Edit canonical lessons, then run npm run generate:books.
% canonical-source-hash: fnv1a64:dc0a7075aef6083a
% canonical-lessons: TA-C329-aracavai, TA-C329-makutam, TA-C329-muticuttu, TA-C329-atci, TA-C329-tutar

\chapter{The Royal Court}
\label{ch:ta-the-royal-court}
\hlchaptermodality{329}

\begin{chapteropening}
\textbf{By the end of this chapter:} \emph{I can describe a royal court and a coronation.}

Complete the last lesson of chapter 329: I can describe a royal court and a coronation.
\end{chapteropening}

\section[\texorpdfstring{\ta{அரசவை} (aracavai)}{அரசவை (aracavai)}]{\ta{அரசவை} (aracavai) — a royal court}

\label{lesson:TA-C329-aracavai}



\begin{quote}
Before the new one: say the Tamil for a princess, then the Tamil for a throne.
\end{quote}

\subsection*{You'll want to know: \ta{அரசவை}}
\textbf{\ta{அரசவை}} — \emph{aracavai} — \textquotedblleft{}a royal court\textquotedblright{}.

\begin{quote}
\ta{புலவர்} \ta{அரசவையில்} \ta{ஒரு} \ta{பாடல்} \ta{பாடினார்}.
\end{quote}

\begin{quote}
\emph{pulavar aracavaiyil oru pāṭal pāṭiṉār.}
\end{quote}

\begin{quote}
The poet sang a song at the royal court.
\end{quote}

\ta{அரசு}, king + \ta{அவை}, assembly.

\begin{cousinweb}
\textbf{\ta{அரசவை}} (\emph{aracavai}) joins \ta{அரசு}, \textquotedblleft{}king, kingdom\textquotedblright{}, which DEDR 201 traces to Sanskrit \emph{rājan}, with \ta{அவை}, \textquotedblleft{}assembly\textquotedblright{}, usually traced to Sanskrit \emph{sabhā}, \textquotedblleft{}assembly, hall\textquotedblright{} (CDIAL 13167).
\end{cousinweb}

\begin{grammarlens}[title={What you've built}]
One more word for telling what happened.
\end{grammarlens}

\subsection*{Your turn}
\begin{itemize}
\raggedright
  \item \emph{Say it:} \emph{aracavai}
  \item \emph{Say it:} \emph{aracavai}, once more
  \item \emph{Say it:} \emph{pulavar aracavaiyil oru pāṭal pāṭiṉār.} — The poet sang a song at the royal court.
  \item \emph{Recall:} say the Tamil for a princess, then the Tamil for a throne, then say \emph{aracavai} again
\end{itemize}

\begin{tcolorbox}[breakable,colback=teal!4,colframe=teal!35!black,title={Before you move on}]
What does \ta{அரசவை} mean? (\textquotedblleft{}A royal court\textquotedblright{}.) Say it once more.
\end{tcolorbox}
\section[\texorpdfstring{\ta{மகுடம்} (makuṭam)}{மகுடம் (makuṭam)}]{\ta{மகுடம்} (makuṭam) — a crown}

\label{lesson:TA-C329-makutam}



\begin{quote}
Before the new one: say the Tamil for a throne, then the Tamil for a royal court.
\end{quote}

\subsection*{You'll want to know: \ta{மகுடம்}}
\textbf{\ta{மகுடம்}} — \emph{makuṭam} — \textquotedblleft{}a crown\textquotedblright{}.

\begin{quote}
\ta{மன்னரின்} \ta{மகுடத்தில்} \ta{பல} \ta{வைரங்கள்} \ta{இருந்தன}.
\end{quote}

\begin{quote}
\emph{maṉṉariṉ makuṭattil pala vairaṅkaḷ iruntaṉa.}
\end{quote}

\begin{quote}
There were many diamonds in the king's crown.
\end{quote}

Also a crowning glory: \ta{வெற்றியின்} \ta{மகுடம்}.

\begin{cousinweb}
\textbf{\ta{மகுடம்}} (\emph{makuṭam}) is borrowed from Sanskrit \emph{mukuṭa}, \emph{makuṭa}, \textquotedblleft{}tiara, crest\textquotedblright{} (CDIAL 10144); CDIAL notes that the Sanskrit word may itself come from a Dravidian word for \textquotedblleft{}top\textquotedblright{}.
\end{cousinweb}

\begin{grammarlens}[title={What you've built}]
One more word for telling what happened.
\end{grammarlens}

\subsection*{Your turn}
\begin{itemize}
\raggedright
  \item \emph{Say it:} \emph{makuṭam}
  \item \emph{Say it:} \emph{makuṭam}, once more
  \item \emph{Say it:} \emph{maṉṉariṉ makuṭattil pala vairaṅkaḷ iruntaṉa.} — There were many diamonds in the king's crown.
  \item \emph{Recall:} say the Tamil for a throne, then the Tamil for a royal court, then say \emph{makuṭam} again
\end{itemize}

\begin{tcolorbox}[breakable,colback=teal!4,colframe=teal!35!black,title={Before you move on}]
What does \ta{மகுடம்} mean? (\textquotedblleft{}A crown\textquotedblright{}.) Say it once more.
\end{tcolorbox}
\section[\texorpdfstring{\ta{முடிசூட்டு} (muṭicūṭṭu)}{முடிசூட்டு (muṭicūṭṭu)}]{\ta{முடிசூட்டு} (muṭicūṭṭu) — to crown a king}

\label{lesson:TA-C329-muticuttu}



\begin{quote}
Before the new one: say the Tamil for a royal court, then the Tamil for a crown.
\end{quote}

\subsection*{You'll want to know: \ta{முடிசூட்டு}}
\textbf{\ta{முடிசூட்டு}} — \emph{muṭicūṭṭu} — \textquotedblleft{}to crown a king\textquotedblright{}.

\begin{quote}
\ta{மக்கள்} \ta{முன்னிலையில்} \ta{இளவரசனுக்கு} \ta{முடிசூட்டினார்கள்}.
\end{quote}

\begin{quote}
\emph{makkaḷ muṉṉilaiyil iḷavaracaṉukku muṭicūṭṭiṉārkaḷ.}
\end{quote}

\begin{quote}
They crowned the prince in front of the people.
\end{quote}

Past (I …): \ta{முடிசூட்டினேன்}, \emph{muṭicūṭṭiṉēṉ}.

The one crowned takes -\ta{க்கு}; \ta{முடி}, crown + \ta{சூட்டு}, to put on.

\begin{cousinweb}
\textbf{\ta{முடிசூட்டு}} (\emph{muṭicūṭṭu}) is a Tamil compound built from \ta{முடி}, \textquotedblleft{}crown, head\textquotedblright{}, and \ta{சூட்டு}, \textquotedblleft{}to adorn, crown, invest\textquotedblright{} (DEDR 2721).
\end{cousinweb}

\begin{grammarlens}[title={What you've built}]
One more word for telling what happened.
\end{grammarlens}

\subsection*{Your turn}
\begin{itemize}
\raggedright
  \item \emph{Say it:} \emph{muṭicūṭṭu}
  \item \emph{Say it:} \emph{muṭicūṭṭu}, once more
  \item \emph{Say it:} \emph{makkaḷ muṉṉilaiyil iḷavaracaṉukku muṭicūṭṭiṉārkaḷ.} — They crowned the prince in front of the people.
  \item \emph{Recall:} say the Tamil for a royal court, then the Tamil for a crown, then say \emph{muṭicūṭṭu} again
\end{itemize}

\begin{tcolorbox}[breakable,colback=teal!4,colframe=teal!35!black,title={Before you move on}]
What does \ta{முடிசூட்டு} mean? (\textquotedblleft{}To crown a king\textquotedblright{}.) Say it once more.
\end{tcolorbox}
\section[\texorpdfstring{\ta{ஆட்சி} (āṭci)}{ஆட்சி (āṭci)}]{\ta{ஆட்சி} (āṭci) — a reign; rule}

\label{lesson:TA-C329-atci}



\begin{quote}
Before the new one: say the Tamil for a crown, then the Tamil for to crown a king.
\end{quote}

\subsection*{You'll want to know: \ta{ஆட்சி}}
\textbf{\ta{ஆட்சி}} — \emph{āṭci} — \textquotedblleft{}a reign; rule\textquotedblright{}.

\begin{quote}
\ta{அவருடைய} \ta{ஆட்சியில்} \ta{நாடு} \ta{அமைதியாக} \ta{இருந்தது}.
\end{quote}

\begin{quote}
\emph{avaruṭaiya āṭciyil nāṭu amaitiyāka iruntatu.}
\end{quote}

\begin{quote}
Under his reign the country was peaceful.
\end{quote}

Also government today: \ta{ஆட்சிக்கு} \ta{வா}, to come to power.

\begin{cousinweb}
\textbf{\ta{ஆட்சி}} (\emph{āṭci}) is native Dravidian: DEDR 5157 gives it under \ta{ஆள்}, \textquotedblleft{}to rule, reign over\textquotedblright{}.
\end{cousinweb}

\begin{grammarlens}[title={What you've built}]
One more word for telling what happened.
\end{grammarlens}

\subsection*{Your turn}
\begin{itemize}
\raggedright
  \item \emph{Say it:} \emph{āṭci}
  \item \emph{Say it:} \emph{āṭci}, once more
  \item \emph{Say it:} \emph{avaruṭaiya āṭciyil nāṭu amaitiyāka iruntatu.} — Under his reign the country was peaceful.
  \item \emph{Recall:} say the Tamil for a crown, then the Tamil for to crown a king, then say \emph{āṭci} again
\end{itemize}

\begin{tcolorbox}[breakable,colback=teal!4,colframe=teal!35!black,title={Before you move on}]
What does \ta{ஆட்சி} mean? (\textquotedblleft{}A reign; rule\textquotedblright{}.) Say it once more.
\end{tcolorbox}
\section[\texorpdfstring{\ta{தூதர்} (tūtar)}{தூதர் (tūtar)}]{\ta{தூதர்} (tūtar) — an envoy; an ambassador}

\label{lesson:TA-C329-tutar}



\begin{quote}
Before the new one: say the Tamil for to crown a king, then the Tamil for a reign.
\end{quote}

\subsection*{You'll want to know: \ta{தூதர்}}
\textbf{\ta{தூதர்}} — \emph{tūtar} — \textquotedblleft{}an envoy; an ambassador\textquotedblright{}.

\begin{quote}
\ta{மன்னர்} \ta{பக்கத்து} \ta{நாட்டுக்கு} \ta{ஒரு} \ta{தூதரை} \ta{அனுப்பினார்}.
\end{quote}

\begin{quote}
\emph{maṉṉar pakkattu nāṭṭukku oru tūtarai aṉuppiṉār.}
\end{quote}

\begin{quote}
The king sent an envoy to the neighbouring country.
\end{quote}

Today also an ambassador.

\begin{cousinweb}
\textbf{\ta{தூதர்}} (\emph{tūtar}) is built on Sanskrit \emph{dūta}, \textquotedblleft{}messenger\textquotedblright{} (CDIAL 6488), with the Tamil -\ta{அர்} of persons.
\end{cousinweb}

\begin{grammarlens}[title={What you've built}]
One more word for telling what happened.
\end{grammarlens}

\subsection*{Your turn}
\begin{itemize}
\raggedright
  \item \emph{Say it:} \emph{tūtar}
  \item \emph{Say it:} \emph{tūtar}, once more
  \item \emph{Say it:} \emph{maṉṉar pakkattu nāṭṭukku oru tūtarai aṉuppiṉār.} — The king sent an envoy to the neighbouring country.
  \item \emph{Recall:} say the Tamil for to crown a king, then the Tamil for a reign, then say \emph{tūtar} again
\end{itemize}

\begin{tcolorbox}[breakable,colback=teal!4,colframe=teal!35!black,title={Before you move on}]
What does \ta{தூதர்} mean? (\textquotedblleft{}An envoy; an ambassador\textquotedblright{}.) Say it once more.
\end{tcolorbox}
